\begin{theorem}[Triangle Inequality]
Let \(V\) be an inner‑product space over \(\mathbb{R}\) or \(\mathbb{C}\) with inner product \(\langle\cdot,\cdot\rangle\) and norm \(\|v\|:=\sqrt{\langle v,v\rangle}\).
For all \(x,y\in V\) one has  
\[
\|x+y\|\le \|x\|+\|y\|.
\]
\end{theorem}

\begin{proof}
\textbf{Degenerate cases.}  
If \(x=0\) then \(\|x+y\|=\|y\|=\|0\|+\|y\|\).  
If \(y=0\) the same argument gives \(\|x+0\|=\|x\|=\|x\|+\|0\|\).  
Hence the inequality holds (with equality) when either vector is the zero vector.
In the sequel we assume \(\,x\neq0\) and \(y\neq0\).

\medskip
\textbf{Squaring the inequality.}  
Since norms are non‑negative, \(\|x+y\|\le\|x\|+\|y\|\) is equivalent to  
\[
\|x+y\|^{2}\le(\|x\|+\|y\|)^{2}. \tag{1}
\]

Using the definition of the norm we expand each side.

\noindent\textit{Left‑hand side:}
\begin{align*}
\|x+y\|^{2}
   &=\langle x+y,\,x+y\rangle\\
   &=\langle x,x\rangle+\langle x,y\rangle+\langle y,x\rangle+\langle y,y\rangle\\
   &=\|x\|^{2}+2\operatorname{Re}\langle x,y\rangle+\|y\|^{2}. \tag{2}
\end{align*}
(The equality \(\langle x,y\rangle+\langle y,x\rangle
   =2\operatorname{Re}\langle x,y\rangle\) follows from conjugate‑symmetry.)

\noindent\textit{Right‑hand side:}
\begin{align}
(\|x\|+\|y\|)^{2}= \|x\|^{2}+2\|x\|\,\|y\|+\|y\|^{2}. \tag{3}
\end{align}

\medskip
\textbf{Application of Cauchy–Schwarz.}  
For all \(u,v\in V\),
\[
|\langle u,v\rangle|\le\|u\|\,\|v\|,
\]
with equality iff \(u\) and \(v\) are linearly dependent.
Consequently,
\[
\operatorname{Re}\langle x,y\rangle\le |\langle x,y\rangle|
                                 \le \|x\|\,\|y\|. \tag{4}
\]

\medskip
\textbf{Deriving the squared inequality.}  
Insert the bound (4) into (2):
\begin{align*}
\|x+y\|^{2}
   &=\|x\|^{2}+2\operatorname{Re}\langle x,y\rangle+\|y\|^{2}\\
   &\le\|x\|^{2}+2\|x\|\,\|y\|+\|y\|^{2}\qquad\text{by (4)}\\
   &= (\|x\|+\|y\|)^{2} \qquad\text{by (3)}.
\end{align*}
Thus (1) holds:
\[
\|x+y\|^{2}\le(\|x\|+\|y\|)^{2}.
\]

\medskip
\textbf{Taking square roots.}  
Both sides of the previous inequality are non‑negative, so taking the
(non‑negative) square root yields
\[
\|x+y\|\le\|x\|+\|y\|.
\]

\medskip
\textbf{Equality case.}  
Equality can occur only if every inequality used above is an equality.
The algebraic expansions (2)–(3) are identities; the potentially strict step
is the Cauchy–Schwarz estimate (4). Hence equality requires
\[
|\langle x,y\rangle|=\|x\|\,\|y\|\quad\text{and}\quad
\operatorname{Re}\langle x,y\rangle=|\langle x,y\rangle|.
\]

The first condition is the equality case of Cauchy–Schwarz, i.e. \(x\) and
\(y\) are linearly dependent: \(y=\lambda x\) for some scalar \(\lambda\) (or one
vector is zero). The second condition forces \(\langle x,y\rangle\) to be a
non‑negative real number, which implies that \(\lambda\) is a non‑negative real
scalar.  

Thus equality holds \emph{iff}
\begin{itemize}
\item \(x=0\) or \(y=0\), or
\item there exists \(\lambda\ge0\) such that \(y=\lambda x\) (the vectors are
positively collinear).
\end{itemize}

Therefore, for any vectors \(x,y\in V\) the triangle inequality
\(\|x+y\|\le\|x\|+\|y\|\) holds, with equality precisely in the cases
described above. \(\square\)
\end{proof}